% Printed 438.
\par\clearpage
\noindent\makebox[0pt][l]{\fontsize{13}{16}\selectfont 438}\hfill{\small\textbf{HÖHENMESSUNGEN.}}\hfill\null\par\medskip
\begin{center}
\setlength{\tabcolsep}{4pt}
\begin{tabular}{l|rrr|rrrr}
&\multicolumn{3}{c|}{Breite}&\multicolumn{4}{c}{\shortstack{[Östl.] Länge\\von Göttingen}}\\
Hannover, Ägydiusthurm & \(52^\circ\)&\(22'\)&\(16^{\prime\prime}{,}4\)&\(-\)&\(0^\circ\)&\(12'\)&\(13^{\prime\prime}{,}8\)\\
\quad---\quad Neustädterthurm &52&22&\(22{,}6\)&\(-\)&0&12&\(52{,}8\)\\
\quad---\quad Marktthurm &52&22&\(24{,}8\)&\(-\)&0&12&\(28{,}4\)\\
\quad---\quad Kreuzthurm &52&22&\(30{,}7\)&\(-\)&0&12&\(37{,}8\)\\
Neustadt am Rübenberge &52&30&\(21{,}8\)&\(-\)&0&28&\(53{,}7\)\\
Celle, südlicher Schlossthurm &52&37&\(31{,}4\)&\(+\)&0&8&\(4{,}9\)\\
\quad---\quad nördlicher Schlossthurm &52&37&\(32{,}9\)&\(+\)&0&8&\(4{,}0\)\\
\quad---\quad Thurm der Stadtkirche &52&37&\(34{,}2\)&\(+\)&0&8&\(16{,}6\)\\
Lüneburg, Johannisthurm &53&14&\(59{,}2\)&\(+\)&0&28&\(11{,}7\)\\
\quad---\quad Michaelisthurm &53&15&\(5{,}5\)&\(+\)&0&27&\(29{,}5\)\\
Hamburg, Michaelisthurm &53&33&\(1{,}8\)&\(+\)&0&2&\(3{,}0\)
\end{tabular}
\end{center}
\par\noindent\hspace*{\parindent}Die Namen der drei Hildesheimschen Thürme kann ich in diesem Augenblick noch nicht angeben; der erste ist ein sogenannter Dachreiter, der zweite
eine Laterne, der dritte nadelförmig. Alle diese Bestimmungen gehen von
dem Platze des \textsc{Reichenbach}schen Meridiankreises in der Göttinger Sternwarte
aus, und zur Rechnung sind die von \textsc{Walbeck} gefundenen Dimensionen des
Erdsphäroids zum Grunde gelegt. Bei der Breite des Brockenhauses weicht
von \textsc{Zachs} astronomische Bestimmung um \(10''\) ab (sein Beobachtungsplatz war
noch etwas südlich vom Thurme); ob dies bloss in der Unvollkommenheit des
von ihm gebrauchten \textsc{Lenoir}schen Kreises seinen Grund habe, wie dieser geschickte Beobachter glaubt, lasse ich dahin gestellt sein; ich möchte aber
doch fast glauben, dass die grosse Menge von Harzgebirgen, welche noch
südlich vom Brocken liegt, während im Norden sogleich das flache Land anfängt, einigen Antheil daran hat, und es wäre gewiss interessant, wenn in Zukunft neue astronomische Beobachtungen mit einem vollkommnern Instrument
auf dem Brocken oder nördlich am Fuss desselben angestellt würden.
\par\medskip\noindent\hspace*{\parindent}Da ich an jedem meiner Dreieckspunkte die Zenithdistanzen aller andern
damit verbundenen mit einem 12-zölligen Multiplicationskreise gemessen habe,
so ergeben sich daraus ihre relativen Höhen mit grosser Genauigkeit; ich
wünsche jedoch, ehe ich diese bekannt mache, erst durch eine zuverlässigere

% Printed 439.
\par\clearpage
\noindent\null\hfill{\small\textbf{HÖHENMESSUNGEN.}}\hfill\makebox[0pt][r]{\fontsize{13}{16}\selectfont 439}\par\medskip
\noindent Verbindung mit der Meeresfläche als bisher stattgefunden hat, die absoluten
Höhen über dieser zu erhalten. Dagegen theile ich Ihnen hier alle meine
bisher erhaltenen Resultate über die terrestrische Refraction mit in folgendem
Tableau, wobei ich nur bemerke, dass ich unter der Refraction die Verschiedenheit der Richtungen des Lichtstrahls an seinen beiden Endpunkten
verstehe. Was einige Astronomen Refraction nennen, ist eigentlich nur die
halbe Refraction.
\begin{center}
\renewcommand{\arraystretch}{1.16}
\begin{tabular}{|l|r|r|r|}\hline
\multicolumn{1}{|c|}{Endpunkte}&\multicolumn{1}{c|}{\shortstack{Krümmung\\des terrestri-\\schen Bogens}}&\multicolumn{1}{c|}{\shortstack{Beobachtete\\Refraction}}&\multicolumn{1}{c|}{Verhältniss}\\\hline
Lichtenberg-Falkenberg &\(2767^{\prime\prime}{,}2\)&\(+345^{\prime\prime}{,}3\)&\(+0{,}1248\)\\
Deister-Falkenberg &\(2282{,}7\)&\(330{,}5\)&\(0{,}1448\)\\
Hohehagen-Brocken &\(2235{,}3\)&\(325{,}6\)&\(0{,}1457\)\\
Lichtenberg-Garssen &\(1961{,}7\)&\(252{,}0\)&\(0{,}1285\)\\
Deister-Garssen &\(1950{,}4\)&\(240{,}5\)&\(0{,}1233\)\\
Hils-Brocken &\(1779{,}1\)&\(292{,}6\)&\(0{,}1645\)\\
Lichtenberg-Deister &\(1597{,}2\)&\(252{,}6\)&\(0{,}1581\)\\
Hohehagen-Hils &\(1529{,}3\)&\(243{,}5\)&\(0{,}1592\)\\
Breithorn-Wilsede &\(1409{,}5\)&\(153{,}6\)&\(0{,}1090\)\\
Brocken-Lichtenberg &\(1372{,}5\)&\(185{,}2\)&\(0{,}1350\)\\
Hils-Deister &\(1316{,}8\)&\(175{,}3\)&\(0{,}1331\)\\
Hils-Lichtenberg &\(1291{,}1\)&\(187{,}1\)&\(0{,}1450\)\\
Hauselberg-Wilsede &\(1232{,}5\)&\(137{,}2\)&\(0{,}1113\)\\
Meridianzeichen-Hils &\(1188{,}9\)&\(187{,}0\)&\(0{,}1573\)\\
Falkenberg-Wilsede &\(1168{,}0\)&\(150{,}7\)&\(0{,}1290\)\\
Falkenberg-Wulfsode &\(1139{,}9\)&\(132{,}1\)&\(0{,}1159\)\\
Falkenberg-Scharnhorst &\(958{,}3\)&\(122{,}4\)&\(0{,}1277\)\\
Garssen-Falkenberg &\(910{,}0\)&\(78{,}3\)&\(0{,}0861\)\\
Falkenberg-Breithorn &\(884{,}0\)&\(109{,}2\)&\(0{,}1235\)\\
Hauselberg-Wulfsode &\(770{,}1\)&\(56{,}1\)&\(0{,}0728\)\\
Wulfsode-Wilsede &\(738{,}9\)&\(+84{,}6\)&\(+0{,}1146\)
\end{tabular}
\end{center}

% Printed 440.
\par\clearpage
\noindent\makebox[0pt][l]{\fontsize{13}{16}\selectfont 440}\hfill{\small\textbf{BRIEFWECHSEL.}}\hfill\null\par\medskip
\begin{center}
\renewcommand{\arraystretch}{1.16}
\begin{tabular}{|l|r|r|r|}\hline
\multicolumn{1}{|c|}{Endpunkte}&\multicolumn{1}{c|}{\shortstack{Krümmung\\des terrestri-\\schen Bogens}}&\multicolumn{1}{c|}{\shortstack{Beobachtete\\Refraction}}&\multicolumn{1}{c|}{Verhältniss}\\\hline
Falkenberg-Hauselberg &\(691^{\prime\prime}{,}7\)&\(+87^{\prime\prime}{,}4\)&\(+0{,}1264\)\\
Meridianzeichen-Hohehagen &\(538{,}6\)&\(111{,}9\)&\(0{,}2078\)\\
Sternwarte-Hohehagen &\(447{,}0\)&\(54{,}3\)&\(0{,}1215\)\\
Scharnhorst-Breithorn &\(363{,}0\)&\(11{,}2\)&\(0{,}0309\)\\
Garssen-Scharnhorst &\(343{,}9\)&\(+20{,}6\)&\(+0{,}0599\)\\
Breithorn-Hauselberg &\(203{,}6\)&\(-23{,}2\)&\(-0{,}1141\)\\
Sternwarte-Meridianzeichen &\(162{,}3\)&\(+36{,}8\)&\(+0{,}2265\)
\end{tabular}
\end{center}
\par\noindent\hspace*{\parindent}Das Mittel aus allen Bestimmungen, mit Rücksicht auf die Länge der
Linien, wäre \(0{,}1306\). Man sieht, dass die Anomalien bei kleinen Entfernungen
viel grösser sind, als bei grossen.
\par\noindent\hspace*{\parindent}An sonnigen Sommertagen ist nach meinen Erfahrungen in den Vormittags- und frühern Nachmittagsstunden in flachen Gegenden die Refraction
des Lichts, so lange es nahe über der Erdfläche wegstreicht, gewöhnlich
negativ, und die Luft ist dann immer so stark undulirend, dass sich keine sehr
scharfe Messungen machen lassen. Das Heliotroplicht ist dann öfters eine
cometenartige Scheibe, zuweilen wohl von einer Minute im Durchmesser, die
sich in den Nachmittagsstunden, so wie die Luft nach und nach günstiger
wird, immer mehr concentrirt und zuletzt in ein feines Fixsternlicht von
grosser Intensität übergeht.
\begin{center}
[2.]\par\medskip
[Der Refractionscoefficient aus den Höhenmessungen von 1824 abgeleitet.]\par\bigskip
\textsc{Gauss} an \textsc{Schumacher}.\quad Göttingen, 28.~November 1824.
\end{center}
\par\noindent\hspace*{\parindent}Erst heute bin ich im Stande, Ihnen das Resultat, welches meine diesjährigen Messungen für die terrestrische Refraction geben, mitzutheilen. Um
die Amplituden genau zu haben, musste ich erst die Längen und Breiten
meiner Dreieckspunkte bestimmen, und dazu mussten erst sämmtliche Dreiecke nach vorgängiger Ausgleichung der Winkel berechnet werden.

